% =============================================================================
% 09_math_challenges.tex -- Mathematical Challenges (~8-10 frames)
% =============================================================================

\section{Mathematical Challenges}

% ---- Frame 1: Overview ----
\begin{frame}{Seven + One Challenges}

  From the companion concept note, adapted for glioma:

  \bigskip

  \begin{enumerate}
    \item The ill-posed inverse problem
    \item Constitutive modeling from dMRI
    \item Registration error amplification
    \item Boundary conditions and unknown forcing
    \item Multimodal fusion across scales
    \item Uncertainty quantification
    \item Validation without ground truth
    \item[\alert{8.}] \alert{Pathological tissue and post-surgical
      deformation}
  \end{enumerate}

  \bigskip

  Challenges 1--7 apply to any brain digital twin.\\
  Challenge 8 is specific to neuro-oncology.

\end{frame}

% ---- Frame 2 ----
\begin{frame}{Challenge 1: The Ill-Posed Inverse Problem}

  The map from MRI observables to material parameters $\bftheta(\bfx)$
  is \emph{non-unique}.

  \bigskip

  \begin{itemize}
    \item Different $(\bftheta, \text{BC})$ combinations produce similar
      predicted observables
    \item Without direct displacement measurements (MRE, tagged MRI),
      the inverse problem is severely underdetermined
  \end{itemize}

  \bigskip

  \textbf{Research questions:}
  \begin{itemize}
    \item Identifiability of mechanical parameters from
      sMRI + dMRI + fMRI alone?
    \item Can SC + FC provide sufficient constraints?
    \item \alert{Glioma advantage:} longitudinal deformation data
      $\bfu(\bfx)$ from SyN provides additional constraints on
      $\bftheta$
  \end{itemize}

\end{frame}

% ---- Frame 3 ----
\begin{frame}{Challenge 2: Constitutive Modeling from dMRI}

  \begin{columns}[T]
    \column{0.55\textwidth}
    \begin{itemize}
      \item DTI $\to$ mechanics mapping is a \emph{modeling assumption},
        not measured \cite{Giordano2014}
      \item Crossing fibers: DTI principal eigenvector is misleading
      \item Partial volume contamination at boundaries
      \item Exponential fiber stiffening:
    \end{itemize}

    \[
      \psi_f = \frac{k_1}{2k_2}
        \left[\exp(k_2 \langle I_4 - 1\rangle^2) - 1\right]
    \]

    Effect is small at physiological strains ($< 1\%$),\\
    significant at traumatic strains ($> 10\%$).

    \column{0.4\textwidth}
    \alert{Glioma-specific:}
    \begin{itemize}
      \item Tumor infiltration alters diffusion properties
      \item Peritumoral edema: increased MD, decreased FA
      \item How to model the tumor--brain interface?
      \item Constitutive law in the infiltration zone?
    \end{itemize}
  \end{columns}

\end{frame}

% ---- Frame 4 ----
\begin{frame}{Challenge 3: Registration and Error Amplification}

  Mechanics depends on \emph{spatial derivatives} (strain $=$
  gradient of displacement). Differentiation amplifies errors.

  \bigskip

  \begin{itemize}
    \item Registration between modalities is itself ill-posed
    \item Registration warps $\neq$ physical displacements
    \item Diffeomorphic registration can introduce deformations the
      mechanical model misinterprets
  \end{itemize}

  \bigskip

  \textbf{Research questions:}
  \begin{itemize}
    \item How does registration uncertainty propagate to stress/strain?
    \item ``Mechanical registration error'': predicted mechanics from
      plausible alternative registrations?
  \end{itemize}

  \medskip

  \alert{Glioma -- concrete example:} A mislocated cavity mask
  ($\sim$50\,mm off) caused global T1w correlation to drop from
  $> 0.85$ to 0.66, with displacement artifacts of 2.4\,mm at the
  actual cavity.  All downstream mechanics would have been built on
  corrupted deformation fields.

\end{frame}

% ---- Frame 5 ----
\begin{frame}{Challenge 4: Boundary Conditions}

  \begin{columns}[T]
    \column{0.55\textwidth}
    \textbf{Standard BC surfaces:}
    \begin{itemize}
      \item Skull: fixed, sliding, or frictional?
      \item Ventricles: pressure BC or FSI?
      \item Falx / tentorium: rigid membrane?
    \end{itemize}

    \bigskip

    \textbf{Glioma-specific BC:}
    \begin{itemize}
      \item \alert{Resection cavity}: new free surface
      \item Cavity filled with CSF: pressure BC?
      \item Cavity wall: scarred tissue, different properties?
      \item Tumor--brain interface: traction-free or infiltrated?
    \end{itemize}

    \column{0.4\textwidth}
    % FIGURE SLOT: schematic of BC surfaces
    \begin{center}
      \fbox{\parbox{0.9\textwidth}{\centering\small
        \textit{Figure:} Boundary\\
        condition surfaces\\
        (skull, ventricles,\\
        cavity)\\[2pt]
        \texttt{figures/bc\_schematic.pdf}
      }}
    \end{center}
  \end{columns}

\end{frame}

% ---- Frame 6 ----
\begin{frame}{Challenge 5: Multimodal Fusion Across Scales}

  \begin{tabular}{lll}
    \toprule
    \textbf{Modality} & \textbf{Scale} & \textbf{Timescale} \\
    \midrule
    dMRI & Sub-voxel ($\mu$m fibers) & Snapshot \\
    fMRI & Voxel-level hemodynamics & Seconds \\
    Mechanics & Mesh element ($\sim$mm) &
      ms (waves) to months (remodeling) \\
    \bottomrule
  \end{tabular}

  \bigskip

  \textbf{Research questions:}
  \begin{itemize}
    \item Spatial scale mismatch: voxel dMRI $\to$ element FE mesh?
    \item Can fMRI temporal information (pulsatile components)
      inform mechanical BCs despite timescale gap?
    \item Joint assimilation framework for sMRI + dMRI + fMRI into
      a mechanical model?
    \item \alert{Glioma:} tumor heterogeneity adds yet another scale
      (microscopic infiltration $\to$ macroscopic deformation)
  \end{itemize}

\end{frame}

% ---- Frame 7 ----
\begin{frame}{Challenge 6: Uncertainty Quantification}

  A credible digital twin must quantify uncertainty.
  Bayesian inversion:
  \[
    p(\bftheta \mid \mathbf{y})
      \propto p(\mathbf{y} \mid \bftheta)\,p(\bftheta)
  \]
  is attractive but computationally prohibitive with fine 3D meshes.

  \bigskip

  \textbf{Approaches:}
  \begin{itemize}
    \item Region-wise parameterization (reduce dimensionality)
    \item Laplace approximation, variational inference
    \item Surrogate models (polynomial chaos, neural network)
    \item Sensitivity indices (Sobol, adjoint-based)
  \end{itemize}

  \bigskip

  \alert{Glioma:} the tumor infiltration zone introduces additional
  parametric uncertainty. ``How far does the tumor extend?'' is
  itself uncertain, and this uncertainty propagates into the
  mechanical prediction.

\end{frame}

% ---- Frame 8 ----
\begin{frame}{Challenge 7: Validation}

  Without MRE, how do we validate?

  \bigskip

  \begin{itemize}
    \item Falsifiability: does anisotropic $\neq$ isotropic?
    \item Cross-modal consistency: stress $\sim$ FC patterns?
    \item Literature comparison: stiffness vs.\ published MRE values
  \end{itemize}

  \bigskip

  \alert{Glioma advantage} (unique to this project):
  \begin{itemize}
    \item Pre-op $\to$ post-op displacement field from SyN
      provides \emph{observed} deformation
    \item Can the FE model predict the spatial pattern of
      $|\bfu(\bfx)|$ under surgery-like boundary conditions?
    \item Jacobian map: does the model predict correct
      expansion/contraction pattern?
  \end{itemize}

  \medskip

  \textbf{Sensitivity analysis as a feature:} The probabilistic tumor
  mask enables running the simulation at thresholds $p > 0.3$, $0.5$,
  $0.7$ -- quantifying how segmentation uncertainty propagates into
  the predicted displacement field.

  \medskip

  This is the closest to \alert{ground-truth displacement data}
  available from clinical MRI alone.

\end{frame}

% ---- Frame 9 ----
\begin{frame}{Challenge 8: Pathological Tissue (Glioma-Specific)}

  New mathematical difficulties introduced by pathology:

  \bigskip

  \begin{itemize}
    \item \textbf{Tumor infiltration zone:} gradual transition from
      normal to pathological tissue properties --- not a sharp boundary
    \item \textbf{Mass effect modeling:} tumor growth as a volumetric
      source term; reaction-diffusion PDE coupled to mechanics
      \cite{Hormuth2021}
    \item \textbf{Post-surgical remodeling:} time-dependent material
      property changes (scarring, gliosis)
    \item \textbf{Topological change:} pre-op domain $\neq$ post-op
      domain (cavity creation)
    \item \textbf{Data contamination:} diffusion/functional signals in
      peri-tumoral region are pathologically altered, not just noisy
  \end{itemize}

  \bigskip

  These are not just ``more noise'' --- they represent \emph{qualitative}
  changes in the mathematical structure of the problem.

\end{frame}
